%! Addition and Multiplication Rules — Unit 3, Lesson 3.3 — Warm-Up
% ONE PAGE, blank and key. Three items at 12pt. Each item is a tool the notes pick up again:
%   Item 1 (decimals .6, .5, .3): the arithmetic of problems 1 and 9 (row 1, row 4).
%   Item 2 (fractions 3/10, 2/9): the arithmetic of problems 5 and 7 (row 3's trees).
%   Item 3 (read the Venn, Lesson 3.2): the drum's Venn, the first display in the notes;
%     24 + 20 = 44 > 40 seeds the trap at problem 2.
% THE DATA: 40 tickets; S 24, E 20, both 12, S only 12, E only 8, neither 8.
\documentclass[12pt]{article}
\usepackage{saar-article}
\usepackage{saar-boxes}
\newcommand{\TallMath}[1]{$\displaystyle #1\rule[-1.4em]{0pt}{3.2em}$}

\begin{document}

\pageheader{Unit 3, Lesson 3.3}{Warm-Up}
% NAMESTRIP: no \namedateperiod here. The student writes their name once, on the
% cover this component is stapled behind.

\vspace{0.15in}

\boxguard
\begin{notesbox}{Spiral Review --- The Concert Raffle}
\normalsize
\begin{enumerate}[leftmargin=1.8em, itemsep=20pt, topsep=6pt]

  \item Compute each value as a decimal.

        \vspace{4pt}
        \renewcommand{\arraystretch}{2.4}
        \begin{tabularx}{\linewidth}{@{} Y Y Y @{}}
        \toprule
        $0.6 + 0.5 - 0.3$ & $0.6 \cdot 0.5$ & $1 - 0.2$ \\
        \midrule
        \blank{2.6cm} & \blank{2.6cm} & \blank{2.6cm} \\
        \bottomrule
        \end{tabularx}

  \item Multiply. Write each answer as a simplified fraction.

        \vspace{4pt}
        \renewcommand{\arraystretch}{2.4}
        \begin{tabularx}{\linewidth}{@{} Y Y @{}}
        \toprule
        $\dfrac{3}{10} \cdot \dfrac{3}{10}$ & $\dfrac{3}{10} \cdot \dfrac{2}{9}$ \\
        \midrule
        \blank{2.6cm} & \blank{2.6cm} \\
        \bottomrule
        \end{tabularx}

  \item \emph{Back to Lesson 3.2.} The raffle drum holds $40$ tickets. The Venn diagram counts
        them: $S$ is \emph{sold to a student} and $E$ is \emph{sold early}.

        \vspace{6pt}
        {\centering
        \begin{tikzpicture}[x=1cm, y=1cm]
          \draw[thick, charcoal] (0,0) rectangle (9,3.3);
          \fill[cerulean, fill opacity=0.18] (3.5,1.6) circle[radius=1.3];
          \fill[goldacc, fill opacity=0.25] (5.3,1.6) circle[radius=1.3];
          \draw[thick, cerulean] (3.5,1.6) circle[radius=1.3];
          \draw[thick, goldacc] (5.3,1.6) circle[radius=1.3];
          \node[font=\bfseries] at (1.9,2.85) {$S$};
          \node[font=\bfseries] at (6.9,2.85) {$E$};
          \node[font=\Large\bfseries] at (2.9,1.6) {12};
          \node[font=\Large\bfseries] at (4.4,1.6) {12};
          \node[font=\Large\bfseries] at (5.9,1.6) {8};
          \node[font=\Large\bfseries] at (8.2,0.55) {8};
        \end{tikzpicture}\par}

        \vspace{8pt}
        \renewcommand{\arraystretch}{1.9}
        \begin{tabularx}{\linewidth}{@{} X p{2.6cm} @{}}
        Tickets inside $S$ & \blank{2.2cm} \\
        Tickets inside $E$ & \blank{2.2cm} \\
        Tickets inside both circles & \blank{2.2cm} \\
        $P(\text{neither})$, as a decimal & \blank{2.2cm} \\
        \end{tabularx}

        \vspace{6pt}
        $S$ count plus $E$ count is more than $40$. In one sentence, why?
        \par\writelines{1}

\end{enumerate}
\end{notesbox}

\end{document}
